\documentclass[a4paper]{article}
\title{Tu idzie tytuł}
\author{tutaj imię i nazwisko autora}
\date{\footnotesize październik 2021}


\usepackage{mathtools,amsthm}    % some advanced mathematics notation
\usepackage{graphicx}            % easy inclusion and manipulation of images
\usepackage{microtype}           % just for fun: really hone in with the typography
\usepackage{booktabs}            % nice looking tables that present the content first and foremost
\usepackage{multicol}            % multiple columns in a local context
\usepackage{pifont}
\usepackage[framemethod=tikz]{mdframed}
\usepackage{tcolorbox}
\usepackage{minted}
\usepackage{fancyhdr}
\usepackage[T1]{fontenc}
% \usepackage[polish]{babel} 
\usepackage{polski}
\usepackage[numbers]{natbib}
\usepackage{wallpaper}


\pagestyle{fancy}
\fancyhf{}

% These are called macros.  They aren't actually this simple, but you can think of it like a smart `find and replace' for most of your needs.
% Macros are case-sensitive.   See section 2.

% These two lines tell LaTeX that I want to
% 1) create a new piece of markup called `term' that takes one `thing' and puts single quotes around it
% 2) create a new piece of markup called `pkg' that takes one `thing' and sets it in a sans-serif font
% 3) redefine an existing piece of markup that doesn't take any `things' and just puts the text you see there.
\newcommand\term[1]{`#1'}
\newcommand\pkg[1]{\textsf{#1}}
\newcommand\writeLaTeX{{\sffamily Overleaf}}
% The above macro (\writeLaTeX) is a good example of when macros come in handy.  This online service recently rebranded itself from 'writeLaTeX' to 'Overleaf'.  All references to the company have been replaced with this one change!  In all honesty this is not *unique* to TeX, but it is exceptionally *easy* in TeX.

% The next few macros use moderately advanced TeX.  They've been fine-tuned for appearance according to personal taste; you should never have to worry about this in your documents.  Ignore these macros except for the logical meanings: `control sequence' (formal term for a TeX macro) and `environment'.
\newcommand\embrace[1]{\texttt{\char`\{#1\char`\}}}
\newcommand\env[1]{\embrace{#1}}
\newcommand\cs[1]{\texttt{\string#1}}
\newcommand\csarg[2]{\cs{#1}\embrace{#2}}

\newtcolorbox{info}{colback=white,colframe=white,left=20pt,right=20pt,top=2pt,bottom=2pt,}

\newtcolorbox{warning}{colback=red!5!white,colframe=red!75!black,left=5pt,right=5pt,top=2pt,bottom=2pt,}

\newtcolorbox{expl}{colback=green!5!white,colframe=green!75!black,left=5pt,right=5pt,top=2pt,bottom=2pt,}

\newtcolorbox{tip}{colback=blue!5!white,colframe=blue!75!black,left=5pt,right=5pt,top=2pt,bottom=2pt,}

\usepackage[hidelinks]{hyperref}
\usepackage{xcolor}
\hypersetup{
    colorlinks,
    linkcolor={blue!50!black},
    citecolor={blue!50!black},
    urlcolor={blue!80!black}
}


% Tu możesz edytować nagłówek strony
\lhead{\footnotesize \copyright~Copyright 2021 by Jakub Walczak}
\rhead{Strona \thepage}
\begin{document}
\maketitle
% A to plik graficzny w lewym dolnym rogu na każdej stronie
\LLCornerWallPaper{0.6}{figs/title.png}


\begin{info}
\footnotesize Tutaj zawarty jest wstęp, ewentualnie jakaś dodatkowa informacja.
\end{info}

\tableofcontents

\newpage
\section{Kwestie organizacyjne realizacji przedmiotu}
Tutaj poruszone są kwestie organizacyjne

% Bloki kolorowe
\begin{warning}
\ding{43} W ramkach czerwonych zawarte są informacje wymagające dodatkowej uwagi - aspekty trudniejsze, albo wymagające ostrożności.
\end{warning}

\begin{tip}
\ding{43} W ramkach niebieskich - dodatkowe informacje lub ciekawostki.
\end{tip}

\begin{expl}
\ding{43} Natomiast w zielonych - dodatkowe wskazówki, jeżeli jesteś napradę początkujący w danym temacie i coś może być jeszcze niejasne.

\end{expl}

\newpage
\section{Rozdział 2.}
W tym rozdziale jest przykładowe środowisko prezentacji kodu

\begin{minted}{text}
The following NEW packages will be INSTALLED:

  ca-certificates    conda-forge/win-64::ca-certificates-2021.5.30-h5b45459_0
  openssl            conda-forge/win-64::openssl-1.1.1l-h8ffe710_0
  pip                conda-forge/noarch::pip-21.2.4-pyhd8ed1ab_0
  python             conda-forge/win-64::python-3.9.7-h7840368_1_cpython
  python_abi         conda-forge/win-64::python_abi-3.9-2_cp39
  setuptools         conda-forge/win-64::setuptools-58.0.4-py39hcbf5309_1
  sqlite             conda-forge/win-64::sqlite-3.36.0-h8ffe710_1
  tzdata             conda-forge/noarch::tzdata-2021a-he74cb21_1
  ucrt               conda-forge/win-64::ucrt-10.0.20348.0-h57928b3_0
  vc                 conda-forge/win-64::vc-14.2-hb210afc_5
  vs2015_runtime     conda-forge/win-64::vs2015_runtime-14.29.30037-h902a5da_5
  wheel              conda-forge/noarch::wheel-0.37.0-pyhd8ed1ab_1

\end{minted}


albo kodu z numeracją linijek:

\begin{minted}[linenos]{python}
>>> def add_subtract(a, b):
...     return a+b, a-b
>>> result = add_subtract(10, 5)
>>> result
(15, 5)
>>> isinstance(result, tuple)
True
\end{minted}


\subsection{Podrozdział 1.1}

\subsection{Inny podrozdział}

\paragraph{A tu nowy paragraf} I przykład cytowania \cite{classes}.


\newpage
\subsection{Zadania} Każde zagadnienie powinno miec skojarzone zadania. Zadania znajdują się zawsze na następnej stronie, po użyciu komendy '\\newpage'

\paragraph{Zadanie \thesection.1}

\paragraph{Zadanie \thesection.2}

\paragraph{Zadanie \thesection.3}

\paragraph{Zadanie \thesection.4}

\newpage


\newpage
\bibliographystyle{IEEEtranN}
\bibliography{references}




\begin{minted}[linenos]{python}

\end{minted}


\end{document}